\sectionkicker{Theme synthesis}
\subsection*{横断比較 — multimodalの「統合」は同じ意味ではない}
\addcontentsline{toc}{subsection}{横断比較 — multimodalの「統合」は同じ意味ではない}
H3、GPT-Live、Qwen-Audio/Wanはすべてmultimodalの前進だが、統合している対象は異なる。生成modalities、リアルタイムinteraction、cross-modal reference、service availabilityを分けると7月の変化が見やすい。
\medskip
\begin{center}
\small
\renewcommand{\arraystretch}{1.16}
\begin{tabularx}{\linewidth}{@{}p{0.20\linewidth}X@{}}
\toprule
観察軸 & 7月の一次資料・論文から読めること \\
\midrule
リアルタイムinteraction & GPT-Liveはsimultaneous listening/speakingを前面に出し、live loopの背後でより深い処理をfrontier modelへ委譲する構成が説明された。 \cite{src-3fb14968493d} \\
cross-modal generation/reference & MiniMax H3はimage・video・audio generationとcross-modal context/referenceを一つのmodel storyとして提示し、2K videoやnative stereo soundはvendor claimとして説明された。 \cite{src-d8e13386974c} \\
service availability & Alibaba Model StudioではWan 2.7 reference-to-video snapshotとQwen-Audio 3.0 TTSのinternational availabilityが月内に更新され、用途別variantとregion差が実運用上の条件になった。 \cite{src-eac8dfbeb78c} \\
公開境界 & H3のweight公開は発表時点では今後数日で行う予定とされており、launchとopen-weight availabilityは分離すべきである。managed realtime modelやregion別API availabilityとも同じ『公開』として扱えない。 \cite{src-d8e13386974c,src-3fb14968493d,src-eac8dfbeb78c} \\
\bottomrule
\end{tabularx}
\renewcommand{\arraystretch}{1.0}
\end{center}
\normalsize
\begin{claimboundary}[この比較の境界]
この表は本号で既に検証・参照している一次資料と論文を横断比較の形へ再配置したものである。vendor / project / author が報告した性能値や優位性は独立再現済みの一般則へ変換せず、本文とSource Notesの帰属境界をそのまま維持する。
\end{claimboundary}
